% Options for packages loaded elsewhere
\PassOptionsToPackage{unicode}{hyperref}
\PassOptionsToPackage{hyphens}{url}
\documentclass[
]{article}
\usepackage{xcolor}
\usepackage[margin=1in]{geometry}
\usepackage{amsmath,amssymb}
\setcounter{secnumdepth}{-\maxdimen} % remove section numbering
\usepackage{iftex}
\ifPDFTeX
  \usepackage[T1]{fontenc}
  \usepackage[utf8]{inputenc}
  \usepackage{textcomp} % provide euro and other symbols
\else % if luatex or xetex
  \usepackage{unicode-math} % this also loads fontspec
  \defaultfontfeatures{Scale=MatchLowercase}
  \defaultfontfeatures[\rmfamily]{Ligatures=TeX,Scale=1}
\fi
\usepackage{lmodern}
\ifPDFTeX\else
  % xetex/luatex font selection
\fi
% Use upquote if available, for straight quotes in verbatim environments
\IfFileExists{upquote.sty}{\usepackage{upquote}}{}
\IfFileExists{microtype.sty}{% use microtype if available
  \usepackage[]{microtype}
  \UseMicrotypeSet[protrusion]{basicmath} % disable protrusion for tt fonts
}{}
\makeatletter
\@ifundefined{KOMAClassName}{% if non-KOMA class
  \IfFileExists{parskip.sty}{%
    \usepackage{parskip}
  }{% else
    \setlength{\parindent}{0pt}
    \setlength{\parskip}{6pt plus 2pt minus 1pt}}
}{% if KOMA class
  \KOMAoptions{parskip=half}}
\makeatother
\usepackage{color}
\usepackage{fancyvrb}
\newcommand{\VerbBar}{|}
\newcommand{\VERB}{\Verb[commandchars=\\\{\}]}
\DefineVerbatimEnvironment{Highlighting}{Verbatim}{commandchars=\\\{\}}
% Add ',fontsize=\small' for more characters per line
\usepackage{framed}
\definecolor{shadecolor}{RGB}{248,248,248}
\newenvironment{Shaded}{\begin{snugshade}}{\end{snugshade}}
\newcommand{\AlertTok}[1]{\textcolor[rgb]{0.94,0.16,0.16}{#1}}
\newcommand{\AnnotationTok}[1]{\textcolor[rgb]{0.56,0.35,0.01}{\textbf{\textit{#1}}}}
\newcommand{\AttributeTok}[1]{\textcolor[rgb]{0.13,0.29,0.53}{#1}}
\newcommand{\BaseNTok}[1]{\textcolor[rgb]{0.00,0.00,0.81}{#1}}
\newcommand{\BuiltInTok}[1]{#1}
\newcommand{\CharTok}[1]{\textcolor[rgb]{0.31,0.60,0.02}{#1}}
\newcommand{\CommentTok}[1]{\textcolor[rgb]{0.56,0.35,0.01}{\textit{#1}}}
\newcommand{\CommentVarTok}[1]{\textcolor[rgb]{0.56,0.35,0.01}{\textbf{\textit{#1}}}}
\newcommand{\ConstantTok}[1]{\textcolor[rgb]{0.56,0.35,0.01}{#1}}
\newcommand{\ControlFlowTok}[1]{\textcolor[rgb]{0.13,0.29,0.53}{\textbf{#1}}}
\newcommand{\DataTypeTok}[1]{\textcolor[rgb]{0.13,0.29,0.53}{#1}}
\newcommand{\DecValTok}[1]{\textcolor[rgb]{0.00,0.00,0.81}{#1}}
\newcommand{\DocumentationTok}[1]{\textcolor[rgb]{0.56,0.35,0.01}{\textbf{\textit{#1}}}}
\newcommand{\ErrorTok}[1]{\textcolor[rgb]{0.64,0.00,0.00}{\textbf{#1}}}
\newcommand{\ExtensionTok}[1]{#1}
\newcommand{\FloatTok}[1]{\textcolor[rgb]{0.00,0.00,0.81}{#1}}
\newcommand{\FunctionTok}[1]{\textcolor[rgb]{0.13,0.29,0.53}{\textbf{#1}}}
\newcommand{\ImportTok}[1]{#1}
\newcommand{\InformationTok}[1]{\textcolor[rgb]{0.56,0.35,0.01}{\textbf{\textit{#1}}}}
\newcommand{\KeywordTok}[1]{\textcolor[rgb]{0.13,0.29,0.53}{\textbf{#1}}}
\newcommand{\NormalTok}[1]{#1}
\newcommand{\OperatorTok}[1]{\textcolor[rgb]{0.81,0.36,0.00}{\textbf{#1}}}
\newcommand{\OtherTok}[1]{\textcolor[rgb]{0.56,0.35,0.01}{#1}}
\newcommand{\PreprocessorTok}[1]{\textcolor[rgb]{0.56,0.35,0.01}{\textit{#1}}}
\newcommand{\RegionMarkerTok}[1]{#1}
\newcommand{\SpecialCharTok}[1]{\textcolor[rgb]{0.81,0.36,0.00}{\textbf{#1}}}
\newcommand{\SpecialStringTok}[1]{\textcolor[rgb]{0.31,0.60,0.02}{#1}}
\newcommand{\StringTok}[1]{\textcolor[rgb]{0.31,0.60,0.02}{#1}}
\newcommand{\VariableTok}[1]{\textcolor[rgb]{0.00,0.00,0.00}{#1}}
\newcommand{\VerbatimStringTok}[1]{\textcolor[rgb]{0.31,0.60,0.02}{#1}}
\newcommand{\WarningTok}[1]{\textcolor[rgb]{0.56,0.35,0.01}{\textbf{\textit{#1}}}}
\usepackage{longtable,booktabs,array}
\newcounter{none} % for unnumbered tables
\usepackage{calc} % for calculating minipage widths
% Correct order of tables after \paragraph or \subparagraph
\usepackage{etoolbox}
\makeatletter
\patchcmd\longtable{\par}{\if@noskipsec\mbox{}\fi\par}{}{}
\makeatother
% Allow footnotes in longtable head/foot
\IfFileExists{footnotehyper.sty}{\usepackage{footnotehyper}}{\usepackage{footnote}}
\makesavenoteenv{longtable}
\usepackage{graphicx}
\makeatletter
\newsavebox\pandoc@box
\newcommand*\pandocbounded[1]{% scales image to fit in text height/width
  \sbox\pandoc@box{#1}%
  \Gscale@div\@tempa{\textheight}{\dimexpr\ht\pandoc@box+\dp\pandoc@box\relax}%
  \Gscale@div\@tempb{\linewidth}{\wd\pandoc@box}%
  \ifdim\@tempb\p@<\@tempa\p@\let\@tempa\@tempb\fi% select the smaller of both
  \ifdim\@tempa\p@<\p@\scalebox{\@tempa}{\usebox\pandoc@box}%
  \else\usebox{\pandoc@box}%
  \fi%
}
% Set default figure placement to htbp
\def\fps@figure{htbp}
\makeatother
\setlength{\emergencystretch}{3em} % prevent overfull lines
\providecommand{\tightlist}{%
  \setlength{\itemsep}{0pt}\setlength{\parskip}{0pt}}
\usepackage{bookmark}
\IfFileExists{xurl.sty}{\usepackage{xurl}}{} % add URL line breaks if available
\urlstyle{same}
% fallback for those not using the hyperref driver hyperxmp:
\makeatletter
\@ifundefined{xmpquote}{\newcommand{\xmpquote}[1]{#1}}{}
\makeatother
\hypersetup{
  pdftitle={Лабораторная работа: Импорт и профилирование данных},
  pdfauthor={Шевчук Д.Р., 02271-ДМ},
  hidelinks,
  pdfcreator={LaTeX via pandoc}}

\title{Лабораторная работа: Импорт и профилирование данных}
\author{Шевчук Д.Р., 02271-ДМ}
\date{2026-10-19}

\begin{document}
\maketitle

{
\setcounter{tocdepth}{2}
\tableofcontents
}
\section{Задание 1}\label{ux437ux430ux434ux430ux43dux438ux435-1}

\begin{Shaded}
\begin{Highlighting}[]
\FunctionTok{library}\NormalTok{(tidyverse)}
\end{Highlighting}
\end{Shaded}

\begin{verbatim}
## -- Attaching core tidyverse packages ------------------------ tidyverse 2.0.0 --
## v dplyr     1.2.1     v readr     2.2.0
## v forcats   1.0.1     v stringr   1.6.0
## v ggplot2   4.0.3     v tibble    3.3.1
## v lubridate 1.9.5     v tidyr     1.3.2
## v purrr     1.2.2     
## -- Conflicts ------------------------------------------ tidyverse_conflicts() --
## x dplyr::filter() masks stats::filter()
## x dplyr::lag()    masks stats::lag()
## i Use the conflicted package (<http://conflicted.r-lib.org/>) to force all conflicts to become errors
\end{verbatim}

\begin{Shaded}
\begin{Highlighting}[]
\FunctionTok{library}\NormalTok{(janitor)}
\end{Highlighting}
\end{Shaded}

\begin{verbatim}
## 
## Присоединяю пакет: 'janitor'
## 
## Следующие объекты скрыты от 'package:stats':
## 
##     chisq.test, fisher.test
\end{verbatim}

\begin{Shaded}
\begin{Highlighting}[]
\FunctionTok{library}\NormalTok{(skimr)}
\FunctionTok{library}\NormalTok{(visdat)}
\FunctionTok{library}\NormalTok{(DataExplorer)}
\end{Highlighting}
\end{Shaded}

\begin{Shaded}
\begin{Highlighting}[]
\FunctionTok{read\_lines}\NormalTok{(}\StringTok{"../data/raw/AirQualityUCI.csv"}\NormalTok{, }\AttributeTok{n\_max =} \DecValTok{5}\NormalTok{)}
\end{Highlighting}
\end{Shaded}

\begin{verbatim}
## [1] "Date;Time;CO(GT);PT08.S1(CO);NMHC(GT);C6H6(GT);PT08.S2(NMHC);NOx(GT);PT08.S3(NOx);NO2(GT);PT08.S4(NO2);PT08.S5(O3);T;RH;AH;;"
## [2] "10/03/2004;18.00.00;2,6;1360;150;11,9;1046;166;1056;113;1692;1268;13,6;48,9;0,7578;;"                                        
## [3] "10/03/2004;19.00.00;2;1292;112;9,4;955;103;1174;92;1559;972;13,3;47,7;0,7255;;"                                              
## [4] "10/03/2004;20.00.00;2,2;1402;88;9,0;939;131;1140;114;1555;1074;11,9;54,0;0,7502;;"                                           
## [5] "10/03/2004;21.00.00;2,2;1376;80;9,2;948;172;1092;122;1584;1203;11,0;60,0;0,7867;;"
\end{verbatim}

\begin{Shaded}
\begin{Highlighting}[]
\FunctionTok{read\_lines}\NormalTok{(}\StringTok{"../data/raw/AirQualityUCI.csv"}\NormalTok{, }\AttributeTok{skip =} \DecValTok{9355}\NormalTok{)}
\end{Highlighting}
\end{Shaded}

\begin{verbatim}
##   [1] "04/04/2005;12.00.00;2,4;1142;-200;12,4;1063;293;603;175;1241;1092;26,9;18,3;0,6406;;"
##   [2] "04/04/2005;13.00.00;2,1;1003;-200;9,5;961;235;702;156;1041;770;28,3;13,5;0,5139;;"   
##   [3] "04/04/2005;14.00.00;2,2;1071;-200;11,9;1047;265;654;168;1129;816;28,5;13,1;0,5028;;" 
##   [4] ";;;;;;;;;;;;;;;;"                                                                    
##   [5] ";;;;;;;;;;;;;;;;"                                                                    
##   [6] ";;;;;;;;;;;;;;;;"                                                                    
##   [7] ";;;;;;;;;;;;;;;;"                                                                    
##   [8] ";;;;;;;;;;;;;;;;"                                                                    
##   [9] ";;;;;;;;;;;;;;;;"                                                                    
##  [10] ";;;;;;;;;;;;;;;;"                                                                    
##  [11] ";;;;;;;;;;;;;;;;"                                                                    
##  [12] ";;;;;;;;;;;;;;;;"                                                                    
##  [13] ";;;;;;;;;;;;;;;;"                                                                    
##  [14] ";;;;;;;;;;;;;;;;"                                                                    
##  [15] ";;;;;;;;;;;;;;;;"                                                                    
##  [16] ";;;;;;;;;;;;;;;;"                                                                    
##  [17] ";;;;;;;;;;;;;;;;"                                                                    
##  [18] ";;;;;;;;;;;;;;;;"                                                                    
##  [19] ";;;;;;;;;;;;;;;;"                                                                    
##  [20] ";;;;;;;;;;;;;;;;"                                                                    
##  [21] ";;;;;;;;;;;;;;;;"                                                                    
##  [22] ";;;;;;;;;;;;;;;;"                                                                    
##  [23] ";;;;;;;;;;;;;;;;"                                                                    
##  [24] ";;;;;;;;;;;;;;;;"                                                                    
##  [25] ";;;;;;;;;;;;;;;;"                                                                    
##  [26] ";;;;;;;;;;;;;;;;"                                                                    
##  [27] ";;;;;;;;;;;;;;;;"                                                                    
##  [28] ";;;;;;;;;;;;;;;;"                                                                    
##  [29] ";;;;;;;;;;;;;;;;"                                                                    
##  [30] ";;;;;;;;;;;;;;;;"                                                                    
##  [31] ";;;;;;;;;;;;;;;;"                                                                    
##  [32] ";;;;;;;;;;;;;;;;"                                                                    
##  [33] ";;;;;;;;;;;;;;;;"                                                                    
##  [34] ";;;;;;;;;;;;;;;;"                                                                    
##  [35] ";;;;;;;;;;;;;;;;"                                                                    
##  [36] ";;;;;;;;;;;;;;;;"                                                                    
##  [37] ";;;;;;;;;;;;;;;;"                                                                    
##  [38] ";;;;;;;;;;;;;;;;"                                                                    
##  [39] ";;;;;;;;;;;;;;;;"                                                                    
##  [40] ";;;;;;;;;;;;;;;;"                                                                    
##  [41] ";;;;;;;;;;;;;;;;"                                                                    
##  [42] ";;;;;;;;;;;;;;;;"                                                                    
##  [43] ";;;;;;;;;;;;;;;;"                                                                    
##  [44] ";;;;;;;;;;;;;;;;"                                                                    
##  [45] ";;;;;;;;;;;;;;;;"                                                                    
##  [46] ";;;;;;;;;;;;;;;;"                                                                    
##  [47] ";;;;;;;;;;;;;;;;"                                                                    
##  [48] ";;;;;;;;;;;;;;;;"                                                                    
##  [49] ";;;;;;;;;;;;;;;;"                                                                    
##  [50] ";;;;;;;;;;;;;;;;"                                                                    
##  [51] ";;;;;;;;;;;;;;;;"                                                                    
##  [52] ";;;;;;;;;;;;;;;;"                                                                    
##  [53] ";;;;;;;;;;;;;;;;"                                                                    
##  [54] ";;;;;;;;;;;;;;;;"                                                                    
##  [55] ";;;;;;;;;;;;;;;;"                                                                    
##  [56] ";;;;;;;;;;;;;;;;"                                                                    
##  [57] ";;;;;;;;;;;;;;;;"                                                                    
##  [58] ";;;;;;;;;;;;;;;;"                                                                    
##  [59] ";;;;;;;;;;;;;;;;"                                                                    
##  [60] ";;;;;;;;;;;;;;;;"                                                                    
##  [61] ";;;;;;;;;;;;;;;;"                                                                    
##  [62] ";;;;;;;;;;;;;;;;"                                                                    
##  [63] ";;;;;;;;;;;;;;;;"                                                                    
##  [64] ";;;;;;;;;;;;;;;;"                                                                    
##  [65] ";;;;;;;;;;;;;;;;"                                                                    
##  [66] ";;;;;;;;;;;;;;;;"                                                                    
##  [67] ";;;;;;;;;;;;;;;;"                                                                    
##  [68] ";;;;;;;;;;;;;;;;"                                                                    
##  [69] ";;;;;;;;;;;;;;;;"                                                                    
##  [70] ";;;;;;;;;;;;;;;;"                                                                    
##  [71] ";;;;;;;;;;;;;;;;"                                                                    
##  [72] ";;;;;;;;;;;;;;;;"                                                                    
##  [73] ";;;;;;;;;;;;;;;;"                                                                    
##  [74] ";;;;;;;;;;;;;;;;"                                                                    
##  [75] ";;;;;;;;;;;;;;;;"                                                                    
##  [76] ";;;;;;;;;;;;;;;;"                                                                    
##  [77] ";;;;;;;;;;;;;;;;"                                                                    
##  [78] ";;;;;;;;;;;;;;;;"                                                                    
##  [79] ";;;;;;;;;;;;;;;;"                                                                    
##  [80] ";;;;;;;;;;;;;;;;"                                                                    
##  [81] ";;;;;;;;;;;;;;;;"                                                                    
##  [82] ";;;;;;;;;;;;;;;;"                                                                    
##  [83] ";;;;;;;;;;;;;;;;"                                                                    
##  [84] ";;;;;;;;;;;;;;;;"                                                                    
##  [85] ";;;;;;;;;;;;;;;;"                                                                    
##  [86] ";;;;;;;;;;;;;;;;"                                                                    
##  [87] ";;;;;;;;;;;;;;;;"                                                                    
##  [88] ";;;;;;;;;;;;;;;;"                                                                    
##  [89] ";;;;;;;;;;;;;;;;"                                                                    
##  [90] ";;;;;;;;;;;;;;;;"                                                                    
##  [91] ";;;;;;;;;;;;;;;;"                                                                    
##  [92] ";;;;;;;;;;;;;;;;"                                                                    
##  [93] ";;;;;;;;;;;;;;;;"                                                                    
##  [94] ";;;;;;;;;;;;;;;;"                                                                    
##  [95] ";;;;;;;;;;;;;;;;"                                                                    
##  [96] ";;;;;;;;;;;;;;;;"                                                                    
##  [97] ";;;;;;;;;;;;;;;;"                                                                    
##  [98] ";;;;;;;;;;;;;;;;"                                                                    
##  [99] ";;;;;;;;;;;;;;;;"                                                                    
## [100] ";;;;;;;;;;;;;;;;"                                                                    
## [101] ";;;;;;;;;;;;;;;;"                                                                    
## [102] ";;;;;;;;;;;;;;;;"                                                                    
## [103] ";;;;;;;;;;;;;;;;"                                                                    
## [104] ";;;;;;;;;;;;;;;;"                                                                    
## [105] ";;;;;;;;;;;;;;;;"                                                                    
## [106] ";;;;;;;;;;;;;;;;"                                                                    
## [107] ";;;;;;;;;;;;;;;;"                                                                    
## [108] ";;;;;;;;;;;;;;;;"                                                                    
## [109] ";;;;;;;;;;;;;;;;"                                                                    
## [110] ";;;;;;;;;;;;;;;;"                                                                    
## [111] ";;;;;;;;;;;;;;;;"                                                                    
## [112] ";;;;;;;;;;;;;;;;"                                                                    
## [113] ";;;;;;;;;;;;;;;;"                                                                    
## [114] ";;;;;;;;;;;;;;;;"                                                                    
## [115] ";;;;;;;;;;;;;;;;"                                                                    
## [116] ";;;;;;;;;;;;;;;;"                                                                    
## [117] ";;;;;;;;;;;;;;;;"
\end{verbatim}

\begin{Shaded}
\begin{Highlighting}[]
\FunctionTok{guess\_encoding}\NormalTok{(}\StringTok{"../data/raw/AirQualityUCI.csv"}\NormalTok{)}
\end{Highlighting}
\end{Shaded}

\begin{verbatim}
## # A tibble: 1 x 2
##   encoding confidence
##   <chr>         <dbl>
## 1 ASCII             1
\end{verbatim}

\section{Задание 2}\label{ux437ux430ux434ux430ux43dux438ux435-2}

\subsection{Параметры
файла}\label{ux43fux430ux440ux430ux43cux435ux442ux440ux44b-ux444ux430ux439ux43bux430}

{\def\LTcaptype{none} % do not increment counter
\begin{longtable}[]{@{}
  >{\raggedright\arraybackslash}p{(\linewidth - 4\tabcolsep) * \real{0.3333}}
  >{\raggedright\arraybackslash}p{(\linewidth - 4\tabcolsep) * \real{0.3333}}
  >{\raggedright\arraybackslash}p{(\linewidth - 4\tabcolsep) * \real{0.3333}}@{}}
\toprule\noalign{}
\begin{minipage}[b]{\linewidth}\raggedright
Параметр
\end{minipage} & \begin{minipage}[b]{\linewidth}\raggedright
Найденное значение
\end{minipage} & \begin{minipage}[b]{\linewidth}\raggedright
Чем подтверждено
\end{minipage} \\
\midrule\noalign{}
\endhead
\bottomrule\noalign{}
\endlastfoot
Разделитель полей & \texttt{;} & В шапке и во всех строках данных поля
отделены именно этим знаком: \texttt{Date;Time;CO(GT);…} \\
Десятичный знак & \texttt{,} & Значения вида \texttt{2,6},
\texttt{11,9}, \texttt{0,7578} в столбцах с дробными числами. Точка в
\texttt{18.00.00} --- не десятичный знак, а разделитель
часов/минут/секунд в поле Time \\
Кодировка & ASCII & Вывод \texttt{guess\_encoding()}:
\texttt{encoding:\ ASCII}. Файл содержит только латиницу, цифры и
базовые знаки препинания, поэтому укладывается в ASCII \\
Служебные строки перед таблицей & 1 (только одна строка заголовка,
данные начинаются со строки 2) & В выводе
\texttt{read\_lines(n\_max\ =\ 5)} строка {[}1{]} --- шапка, строка
{[}2{]} уже первая запись данных (\texttt{10/03/2004;18.00.00;…}) \\
Обозначение пропусков & Код \texttt{-200} & В данных встречается как
заполнитель отсутствующих измерений: \texttt{;-200;} в столбцах
\texttt{NMHC(GT)} и др. \\
Пустые строки в конце файла & Есть, 114 штук & Вывод
\texttt{read\_lines(skip\ =\ 9355)}: после трех строк данных
(12:00--14:00 04/04/2005) идут строки {[}4{]}--{[}117{]} вида
\texttt{";;;;;;;;;"} --- только разделители, без значений \\
Пустые столбцы в конце файла & Есть, 1 последний столбец & Шапка
заканчивается \texttt{…;AH;;}, и каждая строка данных тоже заканчивается
\texttt{;;} --- после столбца \texttt{AH} идет еще один столбец без
имени и без значений \\
\end{longtable}
}

\section{Задание 3}\label{ux437ux430ux434ux430ux43dux438ux435-3}

\subsection{Импорт основного
набора}\label{ux438ux43cux43fux43eux440ux442-ux43eux441ux43dux43eux432ux43dux43eux433ux43e-ux43dux430ux431ux43eux440ux430}

\begin{Shaded}
\begin{Highlighting}[]
\NormalTok{df\_raw }\OtherTok{\textless{}{-}} \FunctionTok{read\_csv2}\NormalTok{(}
  \StringTok{"../data/raw/AirQualityUCI.csv"}\NormalTok{,}
  \AttributeTok{na =} \FunctionTok{c}\NormalTok{(}\StringTok{""}\NormalTok{, }\StringTok{"NA"}\NormalTok{, }\StringTok{"{-}200"}\NormalTok{),}
  \AttributeTok{col\_types =} \FunctionTok{cols}\NormalTok{(}\AttributeTok{Date =} \FunctionTok{col\_character}\NormalTok{(), }\AttributeTok{Time =} \FunctionTok{col\_character}\NormalTok{())}
\NormalTok{)}
\end{Highlighting}
\end{Shaded}

\begin{verbatim}
## i Using "','" as decimal and "'.'" as grouping mark. Use `read_delim()` for more control.
\end{verbatim}

\begin{verbatim}
## New names:
## * `` -> `...16`
## * `` -> `...17`
\end{verbatim}

\begin{Shaded}
\begin{Highlighting}[]
\FunctionTok{glimpse}\NormalTok{(df\_raw) }\CommentTok{\# Первое представление данных (количество строк/столбцов и краткая информация по каждому столбцу)}
\end{Highlighting}
\end{Shaded}

\begin{verbatim}
## Rows: 9,471
## Columns: 17
## $ Date            <chr> "10/03/2004", "10/03/2004", "10/03/2004", "10/03/2004"~
## $ Time            <chr> "18.00.00", "19.00.00", "20.00.00", "21.00.00", "22.00~
## $ `CO(GT)`        <dbl> 2.6, 2.0, 2.2, 2.2, 1.6, 1.2, 1.2, 1.0, 0.9, 0.6, NA, ~
## $ `PT08.S1(CO)`   <dbl> 1360, 1292, 1402, 1376, 1272, 1197, 1185, 1136, 1094, ~
## $ `NMHC(GT)`      <dbl> 150, 112, 88, 80, 51, 38, 31, 31, 24, 19, 14, 8, 16, 2~
## $ `C6H6(GT)`      <dbl> 11.9, 9.4, 9.0, 9.2, 6.5, 4.7, 3.6, 3.3, 2.3, 1.7, 1.3~
## $ `PT08.S2(NMHC)` <dbl> 1046, 955, 939, 948, 836, 750, 690, 672, 609, 561, 527~
## $ `NOx(GT)`       <dbl> 166, 103, 131, 172, 131, 89, 62, 62, 45, NA, 21, 16, 3~
## $ `PT08.S3(NOx)`  <dbl> 1056, 1174, 1140, 1092, 1205, 1337, 1462, 1453, 1579, ~
## $ `NO2(GT)`       <dbl> 113, 92, 114, 122, 116, 96, 77, 76, 60, NA, 34, 28, 48~
## $ `PT08.S4(NO2)`  <dbl> 1692, 1559, 1555, 1584, 1490, 1393, 1333, 1333, 1276, ~
## $ `PT08.S5(O3)`   <dbl> 1268, 972, 1074, 1203, 1110, 949, 733, 730, 620, 501, ~
## $ T               <dbl> 13.6, 13.3, 11.9, 11.0, 11.2, 11.2, 11.3, 10.7, 10.7, ~
## $ RH              <dbl> 48.9, 47.7, 54.0, 60.0, 59.6, 59.2, 56.8, 60.0, 59.7, ~
## $ AH              <dbl> 0.7578, 0.7255, 0.7502, 0.7867, 0.7888, 0.7848, 0.7603~
## $ ...16           <lgl> NA, NA, NA, NA, NA, NA, NA, NA, NA, NA, NA, NA, NA, NA~
## $ ...17           <lgl> NA, NA, NA, NA, NA, NA, NA, NA, NA, NA, NA, NA, NA, NA~
\end{verbatim}

\begin{Shaded}
\begin{Highlighting}[]
\FunctionTok{spec}\NormalTok{(df\_raw) }\CommentTok{\# Подробнее о типах}
\end{Highlighting}
\end{Shaded}

\begin{verbatim}
## cols(
##   Date = col_character(),
##   Time = col_character(),
##   `CO(GT)` = col_double(),
##   `PT08.S1(CO)` = col_double(),
##   `NMHC(GT)` = col_double(),
##   `C6H6(GT)` = col_double(),
##   `PT08.S2(NMHC)` = col_double(),
##   `NOx(GT)` = col_double(),
##   `PT08.S3(NOx)` = col_double(),
##   `NO2(GT)` = col_double(),
##   `PT08.S4(NO2)` = col_double(),
##   `PT08.S5(O3)` = col_double(),
##   T = col_double(),
##   RH = col_double(),
##   AH = col_double(),
##   ...16 = col_logical(),
##   ...17 = col_logical()
## )
\end{verbatim}

\begin{Shaded}
\begin{Highlighting}[]
\FunctionTok{problems}\NormalTok{(df\_raw) }\CommentTok{\# Какие проблемы возникли при считывании}
\end{Highlighting}
\end{Shaded}

\begin{verbatim}
## # A tibble: 0 x 5
## # i 5 variables: row <int>, col <int>, expected <chr>, actual <chr>, file <chr>
\end{verbatim}

\subsection{Пояснение
параметров}\label{ux43fux43eux44fux441ux43dux435ux43dux438ux435-ux43fux430ux440ux430ux43cux435ux442ux440ux43eux432}

{\def\LTcaptype{none} % do not increment counter
\begin{longtable}[]{@{}
  >{\raggedright\arraybackslash}p{(\linewidth - 4\tabcolsep) * \real{0.3333}}
  >{\raggedright\arraybackslash}p{(\linewidth - 4\tabcolsep) * \real{0.3333}}
  >{\raggedright\arraybackslash}p{(\linewidth - 4\tabcolsep) * \real{0.3333}}@{}}
\toprule\noalign{}
\begin{minipage}[b]{\linewidth}\raggedright
Параметр
\end{minipage} & \begin{minipage}[b]{\linewidth}\raggedright
Зачем
\end{minipage} & \begin{minipage}[b]{\linewidth}\raggedright
Что сломалось бы без него
\end{minipage} \\
\midrule\noalign{}
\endhead
\bottomrule\noalign{}
\endlastfoot
\texttt{read\_csv2()} & Уже «знает» два параметра файла: разделитель
\texttt{;} и десятичную запятую (\texttt{sep\ =\ ";"},
\texttt{dec\ =\ ","} --- это его значения по умолчанию) & С
\texttt{read\_csv()} по умолчанию весь файл схлопнулся бы в один
столбец, а числа с запятой (\texttt{2,6}) читались бы как текст \\
\texttt{na\ =\ c("",\ "NA",\ "-200")} & Явный список пропусков: пустые
ячейки (хвостовые строки и два пустых столбца
\texttt{...16}/\texttt{...17}), стандартный \texttt{NA} и код
\texttt{-200} & Без \texttt{-200} датасет выглядел бы «полным»: код
\texttt{-200} остался бы валидным числом и исказил все средние, медианы
и корреляции (это особенность датасета). Без \texttt{""} хвостовые 114
строк-разделителей и пустые столбцы превратились бы не в NA, а в пустые
строки \\
\texttt{col\_types\ =\ cols(Date\ =\ col\_character(),\ Time\ =\ col\_character())}
& Фиксирует оба столбца как текст & Дата \texttt{10/03/2004} в формате
дд/мм/гггг могла бы распарситься неверно или вовсе упасть с ошибкой (в
локали по умолчанию читается мм/дд/гггг --- март и апрель перепутались
бы местами). \texttt{18.00.00} с точкой-разделителем вообще не является
числом или временем --- readr попытался бы угадать и предупредил бы о
проблемах парсинга. Явный тип отключает угадывание для этих столбцов ---
файл гарантированно читается как текст, а собрать \texttt{Date} +
\texttt{Time} можно будет позже \\
\end{longtable}
}

\texttt{problems(df\_raw)} --- Журнал ошибок парсинга --- строки,
которые readr не смог уложить в заявленные типы. Для каждой проблемы
возвращает tibble со столбцами:

\begin{itemize}
\tightlist
\item
  \texttt{row} --- номер строки файла;
\item
  \texttt{col} --- какой столбец сломался;
\item
  \texttt{expected} --- какой тип ожидался;
\item
  \texttt{actual} --- что реально пришло (например, текст \texttt{"NA"}
  вместо числа).
\end{itemize}

В данном случае \texttt{problems(df\_raw)} --- пустой. Это означает, что
все строки файла успешно приведены к заявленным типам столбцов, данные
потеряны не были.

\end{document}
